\label{lecture04}
\subsection*{Цели и маршрут занятия (90 минут)}
В первых трёх лекциях мы изучали, как биты представляют данные и результаты
вычислений. Теперь соединим эти знания с устройством компьютера: кто выполняет
операции, где находятся команды и данные и как они перемещаются между
компонентами. Источник программы занятия --- раздел 2 ПУД курса.

После лекции студент сможет объяснить назначение CPU, RAM, накопителей и GPU;
различить адрес и содержимое памяти; проследить чтение элемента массива;
объяснить роль кэша, контроллера, DMA и прерывания; прочитать короткую
последовательность инструкций NASM. Учебная платформа для конкретных примеров
--- Linux x86-64 с моделью типов LP64. Условные адреса на рисунках не являются
адресами, которые следует подставлять в запускаемую программу.

\begin{center}
\small
\begin{tabular}{|p{0.73\linewidth}|r|}
\hline
Блок & Минуты \\
\hline
Общая картина: инструкции и данные & 0--3 \\
Компоненты компьютера & 3--15 \\
CPU и архитектурные модели & 15--23 \\
Шины и контроллеры & 23--33 \\
Иерархия памяти и кэш & 33--45 \\
Ввод-вывод, DMA, прерывания & 45--54 \\
Адреса, типы, указатели, массивы & 54--62 \\
Сквозной пример C и NASM & 62--76 \\
Стек, куча, время жизни & 76--82 \\
Система целиком & 82--87 \\
Итоги и самопроверка & 87--90 \\
\hline
\end{tabular}
\end{center}

\subsection{Общая картина: два потока}
Чтобы выполнить \texttt{a[1] += 5}, процессору нужны и \emph{инструкция},
описывающая действие, и \emph{данные}, над которыми выполняется действие.
Машинные инструкции тоже представлены байтами. Они не возникают внутри
CPU сами собой: подсистема выборки получает их из памяти.

\begin{quote}
Программа хранится на накопителе; при запуске ОС организует её размещение
в адресном пространстве процесса. Инструкции исполняет процессор, используя
регистры, кэши и оперативную память. Устройства выполняют ввод-вывод.
\end{quote}

Не следует смешивать \textbf{исходный текст} C или NASM, \textbf{исполняемый
файл} с машинными инструкциями и \textbf{процесс} --- выполняющийся экземпляр
программы со своим состоянием. Путь от текста к машинным инструкциям будет
подробно разобран на лекции №7, роль ОС при запуске --- на лекции №5.

\subsection{Компоненты компьютера}
\subsubsection*{CPU и RAM}
\textbf{CPU} (Central Processing Unit, центральный процессор) выполняет
машинные инструкции: арифметику, сравнения, переходы, обращения к памяти.
В одном процессоре может быть несколько \textbf{ядер}. Каждое ядро имеет
средства исполнения инструкций; часть ресурсов, например кэш последнего
уровня, может быть общей. Ядро, логический процессор и физический корпус
процессора --- разные понятия. SMT позволяет одному ядру поддерживать
несколько аппаратных потоков, но не удваивает все исполнительные ресурсы.

\textbf{RAM} (Random Access Memory) --- оперативная память. В обычном ПК это
DRAM: она хранит команды и данные, нужные работающим программам, и требует
питания и периодического обновления заряда. При выключении питания содержимое
обычной RAM теряется. Больший объём RAM позволяет держать больше рабочих
данных, но сам по себе не ускоряет каждое сложение в процессоре.

\textbf{Random access} означает возможность обращаться по адресу, а не
обязательный случайный порядок обращений. Время доступа к реальной DRAM
зависит от состояния контроллера, банков и характера запросов.

\subsubsection*{HDD: механика имеет значение}
\textbf{HDD} хранит данные на вращающихся магнитных пластинах. Для чтения
нужного участка головку требуется переместить к дорожке и дождаться, пока
нужный сектор окажется под ней. Поэтому мелкие случайные обращения дороги:
механическое ожидание измеряется миллисекундами. При последовательном чтении
большого участка эти затраты распределяются на много байтов.

Например, при 7200 оборотах в минуту один оборот занимает
$60/7200\approx0.00833$ с, то есть 8.33 мс. В упрощённой модели среднее
ожидание нужного положения равно половине оборота, около 4.17 мс.
Это \emph{только вращательная задержка}: поиск дорожки, передача и очередь
запросов добавляют время. Контроллер и кэш накопителя также влияют на результат.

\subsubsection*{SSD: флеш-память и контроллер}
\textbf{SSD} обычно использует NAND Flash. Механического перемещения нет,
поэтому случайный доступ значительно быстрее HDD. Однако SSD --- не RAM:
время доступа, интерфейс и правила записи иные. Данные программируются
страницами, а стираются более крупными блоками. Контроллер отображает
логические адреса блоков на физические области Flash, исправляет ошибки
и распределяет износ. Перезапись логического блока не обязательно означает
немедленную перезапись того же физического места.

Типичные характеристики накопителя: ёмкость, задержка, скорость
последовательной передачи и число операций ввода-вывода в секунду (IOPS).
Одно число «MB/s» не описывает все нагрузки. И HDD, и SSD сохраняют данные
без питания, но данные, ещё находящиеся в энергозависимых буферах, могут
требовать отдельного подтверждения сохранения.

\subsubsection*{GPU и видеопамять}
\textbf{GPU} (Graphics Processing Unit) оптимизирован для высокой суммарной
производительности при большом числе подходящих параллельных операций.
Например, над пикселями или элементами матрицы можно выполнять сходные
вычисления. CPU ориентирован в том числе на быстрый ход одного потока
со сложным управлением. Это различие приоритетов, а не запрет CPU на
параллелизм или GPU на условные переходы.

\textbf{Дискретный GPU} обычно имеет собственную видеопамять (VRAM).
Передача данных между RAM и VRAM и запуск работы имеют стоимость: маленькую
задачу может оказаться быстрее выполнить на CPU. \textbf{Встроенный GPU}
обычно использует общую системную RAM, хотя конкретная организация зависит
от платформы. Больший объём VRAM означает возможность хранить больше данных,
а не автоматически большую скорость вычисления.

Путь изображения в упрощённом виде: CPU и драйвер готовят команды и ресурсы;
GPU выполняет работу и формирует изображение в буфере; подсистема вывода
считывает его и передаёт сигнал дисплею. Дисплей не является памятью CPU,
а не вся работа GPU обязательно связана с показом изображения.

\subsubsection*{Плата, периферия, прошивка, питание}
Материнская плата обеспечивает соединения, разъёмы и питание компонентов.
\textbf{Контроллер} управляет обменом с памятью или устройством. В современных
ПК контроллер DRAM и часть линий PCIe часто находятся внутри CPU; чипсет
подключает дополнительные устройства. Конкретная топология различается.

Клавиатура, мышь, сетевой адаптер, накопитель и дисплей взаимодействуют
с системой через свои интерфейсы и контроллеры. \textbf{Драйвер} --- программный
компонент ОС, который знает, как работать с конкретным классом оборудования;
контроллер --- аппаратная часть.

Прошивка платформы, обычно реализующая интерфейсы \textbf{UEFI}, хранится
в энергонезависимой Flash, выполняет начальную инициализацию и помогает
передать управление загрузчику ОС. Исторический термин ROM означает память
только для чтения; современная память прошивки часто допускает обновление.
Блок питания обеспечивает нужные напряжения, охлаждение отводит тепло.
При перегреве процессор может снижать частоту: производительность зависит
и от условий работы, а не только от надписи на корпусе.

\subsection{Что находится внутри ядра CPU}
\begin{description}
\item[Регистры.] Небольшой набор непосредственно доступных инструкциям
хранилищ операндов, адресов и состояния. В x86-64 есть, например,
\texttt{RAX}, \texttt{RBX}, \texttt{RSP}; \texttt{EAX} обозначает младшие
32 бита \texttt{RAX}. Запись в \texttt{EAX} обнуляет старшие 32 бита
\texttt{RAX}. Регистры не являются обычными ячейками RAM с адресами C.
\item[Исполнительные устройства.] АЛУ выполняет целочисленные операции;
есть также устройства работы с плавающей точкой, векторными операциями,
адресами и памятью. Не каждая операция идёт через одно универсальное АЛУ.
\item[Управляющая логика.] Выбирает и декодирует инструкции, организует
их исполнение. Регистр \texttt{RIP} связан с адресом инструкции в
архитектурной модели, а \texttt{RFLAGS} содержит флаги состояния.
\item[Кэши.] Хранят копии инструкций и данных, сокращая обращения к RAM.
Разные уровни кэша имеют разный объём и стоимость доступа.
\end{description}

Учебный цикл: \textbf{выбрать инструкцию --- декодировать --- получить
операнды --- выполнить --- сохранить результат}. Переход меняет дальнейший
порядок исполнения. Это логическая модель: современный CPU использует
конвейер, предсказание переходов и перекрытие работы над несколькими
инструкциями. Одна инструкция не обязана занимать один такт.

При частоте 3 GHz такт длится примерно $1/(3\cdot10^9)$ с, или 0.333 нс.
Частота не равна числу выполненных инструкций в секунду. Важны число
инструкций задачи, степень параллелизма, пропуски кэша и ожидания.

\subsection{Общая и раздельная память команд и данных}
В \textbf{фон-неймановской модели} команды и данные хранятся в общей памяти
и используют общее адресное пространство. В \textbf{гарвардской модели}
память команд и данных, а также пути доступа к ним разделены. Один и тот же
числовой адрес в двух пространствах может обозначать разные места.
В строгой модели обычная команда чтения данных не читает память программ.
Реальные гарвардские процессоры могут предоставлять специальные инструкции
для доступа к коду как к данным.

\begin{center}
\small
\begin{tabular}{|p{0.28\linewidth}|p{0.58\linewidth}|}
\hline
Модель & Организация \\
\hline
Фон Нейман & Общая память команд и данных \\
Гарвард & Раздельные памяти и пути доступа \\
Модифицированная гарвардская & Разделение на части уровней при возможности общей памяти \\
\hline
\end{tabular}
\end{center}

Типичный современный x86-64 сочетает общее адресное пространство с
\emph{раздельными} L1-кэшами инструкций и данных. Поэтому классификация
зависит и от уровня, на котором рассматривается система.

\paragraph{Общая память не означает разрешение исполнять любой байт.}
В простой машине с общей памятью без защиты процессор может пытаться
интерпретировать байты доступной памяти как команды. В современной системе
права страниц ограничивают чтение, запись и исполнение. Например, бит
\texttt{NX/XD} запрещает выборку инструкций из страницы. При попытке
исполнения возникает исключение, а не нормальное выполнение этих байтов.
Стек и куча обычного процесса, как правило, неисполняемы. Точный набор
разрешений определяется архитектурой и ОС. Организация памяти и её защита
--- два разных вопроса.

\subsection{Шины, соединения и контроллеры}
\textbf{Шина} --- средство обмена между компонентами с согласованными
правилами передачи. Классическая учебная схема показывает CPU, RAM и
устройства на общей шине. В современном ПК несколько разных соединений,
в том числе выделенные каналы памяти и коммутируемые последовательные связи.

В логической модели обмена выделяют:
\begin{itemize}
\item \textbf{адрес}: к какому месту или ресурсу обращаемся;
\item \textbf{данные}: что передаём;
\item \textbf{управление}: чтение или запись, готовность, завершение, ошибки.
\end{itemize}
Это не обязательно три отдельных набора проводов: в пакетном соединении
адрес, данные и признаки операции могут передаваться по одним линиям.
Ширина адреса и ширина одновременно передаваемых данных тоже не обязаны
совпадать. Из «64-битный CPU» нельзя вывести, что у каждого его интерфейса
64 физических линии.

\paragraph{Чтение и запись.}
При чтении инициатор формирует запрос с адресом; адресат возвращает данные
и статус. При записи инициатор передаёт адрес и данные, а протокол задаёт,
как сообщить о принятии операции. CPU общается с DRAM через контроллер
памяти: тот переводит запросы в команды конкретных каналов и банков.
Наличие запроса ещё не означает немедленного получения результата.

\paragraph{Современная схема соединений.}
\begin{center}
\small
\begin{tabular}{|p{0.3\linewidth}|p{0.56\linewidth}|}
\hline
Соединение & Типичное назначение \\
\hline
Каналы DDR-памяти & Контроллер памяти CPU и модули DRAM \\
PCIe & GPU, NVMe SSD, сетевые адаптеры; линии и коммутаторы \\
SATA & Подключение HDD и SATA SSD через контроллер \\
USB & Периферия через хост-контроллер и, возможно, концентраторы \\
\hline
\end{tabular}
\end{center}
\textbf{M.2} --- спецификация форм-фактора и разъёма, а не синоним NVMe.
\textbf{NVMe} --- протокол работы с энергонезависимыми накопителями,
обычно через PCIe в локальном ПК. M.2-накопитель может использовать SATA
или PCIe/NVMe; поддержка зависит от устройства и слота. \textbf{PCIe}
не является только «шиной видеокарты»: его используют разные устройства.

\paragraph{Задержка и пропускная способность.}
Задержка --- время ожидания отдельного результата. Пропускная способность
--- объём передачи за единицу времени при заданной нагрузке.
Для грубой модели без перекрытия работы:
\[
T\approx L+\frac{V}{B},
\]
где $L$ --- начальная задержка, $V$ --- объём, $B$ --- скорость передачи.
Например, $L=100$ мкс, $V=1$ MB, $B=1$ GB/s дают около 1.1 мс.
Используем десятичные MB и GB; служебные передачи и очередь не учитываем.
Для маленького запроса доминирует задержка, для большого --- время передачи.
Много параллельных запросов могут повысить суммарную скорость, не уменьшая
задержку каждого. Это полезно помнить при сравнении CPU, GPU и накопителей.
